% !TeX root = ../../Book3.tex
% 纲要标题：### 12.7 标量最大值与 Harnack 方法
% 纲要位置：几何学计划纲要.md，第 2598 行。
% 正文状态：未撰写；本文件目前是框架。

% 纲要内容开始
% 
% * 弱最大值原理及零阶项符号
% * 强最大值原理和屏障
% * Hopf 边界引理的条件
% * Moser 迭代与所选 Harnack 不等式
% * 系统不具备同一结论的原因
% 
% ---
% 
% 纲要内容结束

\section{标量最大值与 Harnack 方法}
\label{g3:sec:1207}

待撰写

\subsection{弱最大值原理及零阶项符号}
\label{g3:subsec:1207:01}
待撰写

\subsection{强最大值原理和屏障}
\label{g3:subsec:1207:02}
待撰写

\subsection{Hopf 边界引理的条件}
\label{g3:subsec:1207:03}
待撰写

\subsection{Moser 迭代与所选 Harnack 不等式}
\label{g3:subsec:1207:04}
待撰写

\subsection{系统不具备同一结论的原因}
\label{g3:subsec:1207:05}
待撰写
% 纲要条目排为编号小节标题，并进入目录与书签；说明段仍保留于纲要注释中。
